\chapter{Explorer}
\label{ch:explorer}
% BUDGET : 7 pages (cf. memoire/PLAN.md, section 4).
% THESE DU CHAPITRE : un classement de politiques publiques n'a de sens que relativement a un
% forcage, et le dispositif qui le produit doit enoncer ce qu'il ne peut PAS trancher.
%
% NOTE : les trois tableaux (grille, front budget, realloc) sont CONSERVES -- ils sont
% entierement bases sur des macros generees, il n'y a rien a retaper.

Le modèle porté et calibré, la commande initiale peut enfin être posée : que produirait telle
politique publique, et que produirait-elle si le contexte se dégradait ? Ce chapitre présente
le dispositif construit pour y répondre, ce qu'il a effectivement mesuré, et — avec autant de
soin — ce qu'il ne permet pas de conclure.

\section{Le dispositif expérimental}
\label{sec:dispositif}

Un scénario, ici, n'est pas un objet mais trois, et ils sont séparés \emph{parce qu'ils n'ont
pas le même auteur}. Une \textbf{politique} est ce qu'une puissance publique décide : un
plafond d'azote, une enveloppe de subventions, un plancher d'emploi, une part d'inertie par
rapport à l'assolement existant. Un \textbf{forçage} est ce que le climat et les marchés
imposent : des prix, des rendements, des coûts d'intrants multipliés par un facteur. Un
\textbf{balayage} fait varier un seul seuil et trace le front des compromis atteignables. Le
catalogue déclare \num{\mesPolitiques} politiques, \num{\mesForcages} forçages et
\num{\mesBalayages} fronts.

La figure~\ref{fig:spectre} situe les uns et les autres. Elle sert à deux choses : montrer
que le catalogue de politiques est ordonné le long d'un \emph{axe} — et non pas assemblé au
hasard —, et montrer d'un coup d'\oe il quelles cases ont effectivement été calculées. Ce
second point est une limite du travail, et il vaut mieux la voir tout de suite que la
découvrir au §~\ref{sec:limites}.

\begin{figure}[tbp]
  \centering
  \begin{tikzpicture}[
      x=1mm, y=1mm, font=\scriptsize,
      >={Stealth[length=1.8mm]},
      axe/.style     = {<->, thick, black!75},
      resolu/.style  = {circle, draw=black!80, fill=black!80, inner sep=1.1pt},
      absent/.style  = {circle, draw=black!55, fill=white, inner sep=1.1pt},
      etiq/.style    = {align=center, text width=24mm, inner sep=1pt},
      borne/.style   = {align=center, text width=24mm, inner sep=1pt, font=\scriptsize\itshape}]

    % ================= LES DIX POLITIQUES =================
    % L'axe va de 12 a 138 et non de 0 a 142 : les etiquettes sont centrees sur leur repere et
    % debordent donc de la moitie de `text width` de chaque cote. Sans cette marge de 12 mm, la
    % figure est plus large que la justification et TeX signale un Overfull hbox de 47 pt.
    \draw[axe] (12,0) -- (138,0);
    \node[anchor=north west, font=\scriptsize\itshape, text=black!60] at (0,-4)
      {moins de contrainte publique};
    \node[anchor=north east, font=\scriptsize\itshape, text=black!60] at (150,-4)
      {davantage de contrainte publique};
    \foreach \x in {12,30,52,70,90,106,122,138} \draw[black!45] (\x,-1.5) -- (\x,1.5);

    \node[resolu] at (12,0)  {}; \node[etiq, above=2mm] at (12,0)  {\textbf{P1}\\dérégulation totale};
    \node[absent] at (30,0)  {}; \node[etiq, above=2mm] at (30,0)  {\textbf{P2}\\accélérationnisme industriel};
    \node[resolu] at (70,0)  {}; \node[etiq, above=2mm] at (70,0)  {\textbf{P4}\\statu quo};
    \node[resolu] at (106,0) {}; \node[etiq, above=2mm] at (106,0) {\textbf{P7}\\Écophyto réglementaire};
    \node[absent] at (138,0) {}; \node[etiq, above=2mm] at (138,0) {\textbf{P10}\\bifurcation agroécologique};

    \node[resolu] at (30,0)  {}; \node[etiq, below=6mm] at (30,0)  {\textbf{P5}\\austérité budgétaire};
    \node[absent] at (52,0)  {}; \node[etiq, below=6mm] at (52,0)  {\textbf{P3}\\intensification modérée};
    \node[resolu] at (90,0)  {}; \node[etiq, below=6mm] at (90,0)  {\textbf{P6}\\verdissement incitatif};
    \node[resolu] at (122,0) {}; \node[etiq, below=6mm] at (122,0) {\textbf{P8}\\transition agroécologique};

    \node[resolu] at (2,-22) {};
    \node[anchor=north west, align=left, text width=142mm] at (5,-20)
      {\textbf{P9} souveraineté alimentaire --- \emph{hors axe} : elle contraint la destination
       de la production, pas son intensité};

    % ================= LES DOUZE FORÇAGES =================
    \draw[axe] (12,-38) -- (138,-38);
    \foreach \x in {12,39,138} \draw[black!45] (\x,-39.5) -- (\x,-36.5);
    \node[absent] at (12,-38)  {};
    \node[resolu] at (39,-38)  {};
    \node[resolu] at (138,-38) {};
    \node[borne, above=2mm] at (12,-38)  {\textbf{F10}\\conjoncture favorable};
    \node[borne, above=2mm] at (39,-38)  {\textbf{F0}\\nominal (témoin nul)};
    \node[borne, above=2mm] at (138,-38) {\textbf{F9}\\crise systémique};

    \draw[black!55] (52,-40) -- (52,-43) -- (126,-43) -- (126,-40);
    \node[align=center, text width=130mm, below=0.5mm] at (75,-43)
      {\textbf{huit forçages à un seul levier} --- climat et sanitaire : F1 cyclone
       majeur$^{\bullet}$, F2 sécheresse sévère, F3 dérive climatique, F4 choc sanitaire ;
       marchés et politique : F5 effondrement des prix d'export, F6 choc sur les
       intrants$^{\bullet}$, F7 inflation alimentaire$^{\bullet}$, F8 retrait du
       \textsc{posei}, F11 artificialisation des terres};
  \end{tikzpicture}
  \caption{Les politiques et les forçages du catalogue, situés les uns par rapport aux autres.
  Un disque plein signale qu'au moins une cellule a été résolue (marque $^{\bullet}$ pour les
  forçages), un cercle vide qu'aucune ne l'a été. La position sur l'axe des politiques est
  \emph{ordinale} : elle traduit la sévérité des contraintes imposées, non une distance
  mesurée. P2 et P5 partagent une abscisse parce qu'aucune des deux n'impose de contrainte
  environnementale ; elles s'opposent sur les aides, dimension que cet axe unique ne porte pas.
  La répartition des cercles vides n'est pas neutre : \textbf{la borne dérégulatrice de l'axe a
  été calculée, sa borne de bifurcation non} (§~\ref{sec:limites}).}
  \label{fig:spectre}
\end{figure}

Cette séparation n'est pas seulement conceptuelle : elle décide de ce que coûte une campagne.
Le produit cartésien complet vaudrait \num{\mesCellulesTotal} cellules, soit environ cinq jours
de calcul continu pour la seule grille, sans les fronts. Un plan explicite désigne donc, pour
chaque politique, les forçages qu'elle traverse et les fronts qu'on trace sous elle ; et
surtout, \textbf{un front n'est jamais croisé avec les forçages de sa cellule}, puisqu'un front
est déjà une résolution par point. \num{\mesCellulesGrille} cellules de grille et
\num{\mesSolvesProspectifs} dossiers de résultats prospectifs au total constituent le matériau
de ce chapitre. Cette grille est \emph{creuse}, et c'est une limite du travail plutôt qu'un
choix de méthode : une cellule coûte une heure.

Trois décisions de méthode méritent d'être défendues. \textbf{Un}, les seuils du catalogue sont
des \emph{fractions choisies de totaux mesurés} sur un run de référence, et non des chiffres
inventés : un plancher d'emploi est un pourcentage d'un maximum mesuré, un plafond d'azote une
fraction de l'azote effectivement épandu. L'annexe~\ref{ann:scenarios} donne, pour chaque
seuil, la fraction \emph{et} le total dont elle est tirée. \textbf{Deux}, les forçages ne sont
\textbf{pas moyennés} : une moyenne affirmerait une distribution de probabilité que personne n'a
choisie. Une politique est donc décrite par le vecteur de ses résultats sous chaque forçage,
avec les critères classiques de décision dans l'incertain — pire cas \parencite{wald1945},
regret maximal \parencite{savage1951}, domaine de viabilité \parencite{starr1963} — et
l'infaisabilité y compte comme le \emph{pire} résultat, non comme une donnée manquante.
\textbf{Trois}, les fronts sont tracés par $\varepsilon$-contrainte \parencite{haimes1971,
mavrotas2009} et non par somme pondérée. La raison est un théorème, pas un banc d'essai : une
somme pondérée ne peut atteindre que les points situés sur l'enveloppe convexe du front, et ses
poids encoderaient une valeur politique déguisée en paramètre technique.

Deux précautions numériques s'y ajoutent, sans lesquelles les chiffres de la suite ne
vaudraient rien, et l'annexe~\ref{ann:scenarios} en donne le détail. La première est
l'amorçage en chaîne : sous une contrainte du type \enquote{au plus $\varepsilon$}, une
solution admissible pour un $\varepsilon$ serré l'est encore pour tout $\varepsilon$ plus
lâche, ce qui permet de ne payer qu'une résolution à froid pour six à chaud — à condition de
résoudre du plus serré au plus lâche, c'est-à-dire dans l'ordre \emph{inverse} de la lecture du
fichier. La seconde est le contrôle de faisabilité en relaxation \gls{lp} — deux minutes contre
une heure —, dont la logique est à sens unique : \emph{infaisable en \gls{lp} implique
infaisable en entiers, la réciproque est fausse}. C'est un filtre, non une garantie, et il a
suffi à écarter une erreur de conception avant qu'aucun lot ne tourne.

\section{Ce que la grille montre}
\label{sec:grille}

Le résultat central de la grille tient en une phrase : \textbf{le classement des politiques
dépend du forçage}. Une politique n'est pas \enquote{bonne} dans l'absolu ; elle est performante
dans un contexte et exposée dans un autre. C'est ce qu'une grille est faite de montrer, et
c'est exactement ce qu'un scénario unique masque.

Le tableau~\ref{tab:grille} donne les cellules résolues. Trois lectures s'en dégagent, dont une
ne tient pas.

\textbf{La première tient, et elle est brutale.} Sous crise systémique, la marge brute du
territoire devient \emph{négative} pour les trois politiques testées. C'est le pire cas au sens
de Wald, et il est atteint par toutes : aucune des politiques instruites n'offre d'abri contre
ce forçage-là.

\textbf{La deuxième tient, et c'est la plus instructive.} La politique de souveraineté
alimentaire P9 est la \emph{meilleure} sous le forçage nominal — \num{\grillePneufFzero} contre
\num{\grillePquatreFzero} millions d'euros de marge brute pour le statu quo — et la
\emph{pire} sous crise systémique, avec \num{\grillePneufFneuf} contre
\num{\grillePquatreFneuf} millions. Le classement s'inverse complètement d'une colonne à
l'autre. Le regret de Savage, calculé contre la meilleure politique \emph{du même forçage} —
\enquote{qu'aurais-je perdu à choisir celle-ci, sachant que ce contexte-là est survenu} — rend
l'inversion visible d'un coup d'\oe il (figure~\ref{fig:regret}) : P9 passe d'un regret nul
sous F0 à \SI{\regretPneufFneuf}{\percent} sous F9, tandis que le statu quo fait exactement
l'inverse.

\textbf{La troisième ne tient pas.} La marge brute \emph{contient} la subvention : classer sur
cette colonne des politiques dont l'objet même est de déplacer les aides, c'est en partie
constater qu'elles subventionnent davantage. La colonne \enquote{hors aide} isole ce que
l'assolement produit par lui-même, et le classement change — sous crise systémique, c'est
alors P8 qui devient la plus exposée (\num{\grilleNetPhuitFneuf} millions) et non P9
(\num{\grilleNetPneufFneuf}). Aucune des deux lectures n'est fausse : elles répondent à deux
questions différentes. Ce qui serait fautif serait d'en choisir une en silence, et c'est
pourquoi le tableau porte les deux.

Reste à nommer ce qui, dans ce modèle, limite réellement les politiques. Ce n'est ni la terre
ni l'azote, c'est le \textbf{travail}. Le plafond de main d'\oe uvre autorise
\SI{6252740}{\hour}, soit \SI{3891}{\etp} au desserrement nominal, et le run de référence en
sature déjà \SI{90}{\percent} ; le desserrer ne suffit d'ailleurs pas, car ce ne sont pas les
heures qui manquent mais \emph{les cultures capables de les absorber} sous les autres plafonds.
Il en découle un résultat de politique publique et non une observation technique :
\textbf{dans ce modèle, une politique environnementale ambitieuse plafonne l'emploi agricole
qu'elle peut soutenir}. On se gardera de l'explication trop rapide — \enquote{le travail et
l'\gls{ift} sont corrélés} est vrai dans la table des itinéraires et ne suffit pas à rendre
compte d'un résultat qui passe par le jeu des contraintes.

Un dernier fait, contre-intuitif : un \emph{forçage} peut annuler une politique. Sous crise
systémique, le plafond d'azote de P8 ne mord jamais, la crise réduisant d'elle-même les surfaces
intensives. La politique environnementale ne coûte alors rien — et n'apporte rien.

\begin{table}[tbp]
  \centering
  \caption{Marge brute par cellule (M\eur). \emph{Brute} inclut la subvention ;
  \emph{hors aide} en déduit le transfert. La dague marque un solve arrêté à la limite de temps,
  donc prouvablement sous-optimal.}
  \label{tab:grille}
  \small
  \begin{tabularx}{\textwidth}{@{}>{\raggedright\arraybackslash}Xrrrrrr@{}}
    \toprule
    & \multicolumn{2}{c}{\textbf{F0 nominal}} & \multicolumn{2}{c}{\textbf{F6 intrants}}
    & \multicolumn{2}{c}{\textbf{F9 crise}} \\
    \cmidrule(lr){2-3}\cmidrule(lr){4-5}\cmidrule(lr){6-7}
    & brute & hors aide & brute & hors aide & brute & hors aide \\
    \midrule
    P4 statu quo & \num{\grillePquatreFzero}\grillePquatreFzeroFlag & \num{\grilleNetPquatreFzero}
                 & \num{\grillePquatreFsix}\grillePquatreFsixFlag   & \num{\grilleNetPquatreFsix}
                 & \num{\grillePquatreFneuf}\grillePquatreFneufFlag & \num{\grilleNetPquatreFneuf} \\
    P6 verdissement & \num{\grillePsixFzero}\grillePsixFzeroFlag & \num{\grilleNetPsixFzero}
                 & \num{\grillePsixFsix}\grillePsixFsixFlag & \num{\grilleNetPsixFsix}
                 & --- & --- \\
    P7 Écophyto  & \num{\grillePseptFzero}\grillePseptFzeroFlag & \num{\grilleNetPseptFzero}
                 & \num{\grillePseptFsix}\grillePseptFsixFlag & \num{\grilleNetPseptFsix}
                 & --- & --- \\
    P8 agroécologie & \num{\grillePhuitFzero}\grillePhuitFzeroFlag & \num{\grilleNetPhuitFzero}
                 & --- & ---
                 & \num{\grillePhuitFneuf}\grillePhuitFneufFlag & \num{\grilleNetPhuitFneuf} \\
    P9 souveraineté & \num{\grillePneufFzero}\grillePneufFzeroFlag & \num{\grilleNetPneufFzero}
                 & --- & ---
                 & \num{\grillePneufFneuf}\grillePneufFneufFlag & \num{\grilleNetPneufFneuf} \\
    \bottomrule
  \end{tabularx}
\end{table}

\begin{figure}[tbp]
  \centering
  \begin{tikzpicture}[x=1.25mm, y=1mm, font=\scriptsize]
    % Regret en % de la meilleure politique du MEME forcage. Barres empilees par politique :
    % F0 en fonce, F6 en gris moyen, F9 en gris clair. Un regret nul est marque par un trait.
    \foreach \g in {0,10,20,30,40,50} {
      \draw[black!12] (\g,2) -- (\g,-46);
      \node[below, inner sep=1.5pt] at (\g,-46) {\g};
    }
    \draw[black!50] (0,2) -- (0,-46);
    % --- P4 statu quo
    \node[left, inner xsep=3pt] at (0,-2) {P4 statu quo};
    \fill[black!70] (0,0)    rectangle (\regretPquatreFzero,-1.6);
    \fill[black!40] (0,-2)   rectangle (\regretPquatreFsix,-3.6);
    \fill[black!15] (0,-4)   rectangle (0.2,-5.6);
    % --- P6 verdissement
    \node[left, inner xsep=3pt] at (0,-11) {P6 verdissement};
    \fill[black!70] (0,-9)   rectangle (\regretPsixFzero,-10.6);
    \fill[black!40] (0,-11)  rectangle (0.2,-12.6);
    % --- P7 Ecophyto
    \node[left, inner xsep=3pt] at (0,-20) {P7 Écophyto};
    \fill[black!70] (0,-18)  rectangle (\regretPseptFzero,-19.6);
    \fill[black!40] (0,-20)  rectangle (\regretPseptFsix,-21.6);
    % --- P8 agroecologie
    \node[left, inner xsep=3pt] at (0,-29) {P8 agroécologie};
    \fill[black!70] (0,-27)  rectangle (\regretPhuitFzero,-28.6);
    \fill[black!15] (0,-29)  rectangle (\regretPhuitFneuf,-30.6);
    % --- P9 souverainete
    \node[left, inner xsep=3pt] at (0,-38) {P9 souveraineté};
    \fill[black!70] (0,-36)  rectangle (0.2,-37.6);
    \fill[black!15] (0,-38)  rectangle (\regretPneufFneuf,-39.6);
    % --- Legende
    \fill[black!70] (0,-52)  rectangle (2.5,-53.6); \node[right] at (2.5,-52.8) {F0 nominal};
    \fill[black!40] (16,-52) rectangle (18.5,-53.6); \node[right] at (18.5,-52.8) {F6 intrants};
    \fill[black!15] (33,-52) rectangle (35.5,-53.6); \node[right] at (35.5,-52.8) {F9 crise};
  \end{tikzpicture}
  \caption{Regret par politique et par forçage, en pourcentage de la meilleure politique du même
  forçage. Une barre nulle (trait) signifie \enquote{c'était le meilleur choix, sachant ce qui
  est arrivé} ; les cellules non résolues n'ont pas de barre. La lecture qui compte est
  \emph{verticale} : P9 est sans regret sous le forçage nominal et le plus regretté sous crise,
  P4 fait l'inverse. Aucune politique n'est courte partout.}
  \label{fig:regret}
\end{figure}

\section{Le coût marginal : euro public et azote}
\label{sec:fronts}

\subsection*{L'enveloppe publique}

Le résultat central de ce chapitre est un front, et il est contre-intuitif : \textbf{la marge
brute varie en sens inverse du plafond de subventions}, de façon monotone sur les sept points
(tableau~\ref{tab:front-budget}). Serrer l'enveloppe publique de
\SI{\frontsubvQuatreSeuilHaut}{\euros} à \SI{\frontsubvQuatreunSeuil}{\euros}
\emph{augmente} la marge brute du territoire de plus de \SI{14}{\mega\euros}, tout en faisant
reculer l'objectif de \num{\frontsubvQuatreObjHaut} à \num{\frontsubvQuatreObjBas} millions.

Les deux mouvements ne se contredisent pas, et leur écart est instructif. L'objectif est une
marge \emph{ajustée du risque} : la pénalité de risque passe de \SI{18}{\percent} de la marge
brute au plafond le plus haut à \SI{34}{\percent} au plus bas. Autrement dit, privé
d'enveloppe, le modèle va chercher des cultures plus rémunératrices et plus exposées.
\textbf{Dans ce modèle, l'aide publique n'achète pas de la marge, elle achète de la stabilité.}

Avant d'en tirer quoi que ce soit, le front se valide par ses points de contrôle. Deux des sept
points portent un plafond supérieur à ce que le statu quo dépense spontanément : ils ne
contraignent donc rien, et l'objectif doit y retomber exactement sur le run non balayé
(\num{\frontsubvQuatreTemoin} millions). C'est le cas, au centime. Un front dont les points de
contrôle ne se comportent pas ainsi est un front mal paramétré — et c'est un défaut qui s'est
effectivement produit lors de cette campagne (§~\ref{sec:limites}).

Le front donne aussi quelque chose que le prix dual ne donne pas : au point le plus serré, le
dual vaut près du double de la pente \emph{moyenne} du front, si bien que l'extrapoler sur tout
l'intervalle surestimerait l'effet d'environ trois quarts. Le front étant convexe, le dual n'est
valable qu'au voisinage de son point.

Ce que l'enveloppe maintient se lit enfin à l'hectare près (tableau~\ref{tab:ann-realloc}).
Entre le statu quo libre et le plafond le plus serré, la banane d'exportation passe de
\SI{2004}{\ha} à \SI{95}{\ha}, soit \SI{-95}{\percent} ; la canne recule de \SI{2231}{\ha} ; le
maraîchage double, les agrumes passent de \SI{9}{\ha} à \SI{657}{\ha}, l'igname double et la
prairie gagne \SI{1539}{\ha}. C'est l'énoncé le plus direct que ce mémoire produise sur la
question du projet : \emph{dans les termes du modèle, ce qui tient la banane et la canne sur la
terre guadeloupéenne, c'est l'enveloppe publique} — et desserrer cette contrainte réoriente
environ \SI{2100}{\ha} vers des cultures vivrières.

Cette phrase appelle immédiatement sa réserve, et il serait malhonnête de la reporter à la fin
du chapitre : rien, dans le modèle, ne demande si les \SI{2237}{\ha} de maraîchage ainsi
produits trouvent un acheteur (§~\ref{sec:limites}).

\begin{table}[tbp]
  \centering
  \caption{Front marge $\times$ enveloppe publique sous P4, en euros. L'astérisque marque les
  points de contrôle : l'objectif y retombe exactement sur le run non balayé
  (\num{\frontsubvQuatreTemoin}~M\eur), donc le plafond n'y contraint rien.}
  \label{tab:front-budget}
  \small
  % Le tabular entier est genere (entete et filets compris) : le \input de LaTeX ouvre une
  % cellule et ne peut donc pas vivre A L'INTERIEUR d'un alignement.
  % Genere par memoire/build_chiffres.py -- ne pas editer.
\begin{tabular}{@{}lrr@{}}
  \toprule
  \textbf{Plafond} & \textbf{Niveau atteint} & \textbf{Objectif (M\eur)} \\
  \midrule
  \num{35900000} & \num{35898271} & \num{74.36} \\
  \num{43000000} & \num{42986537} & \num{76.84} \\
  \num{50300000} & \num{50070701} & \num{77.96} \\
  \num{57400000} & \num{57111266} & \num{79.18} \\
  \num{64600000} & \num{64590701} & \num{80.62} \\
  \num{72000000} & \num{71277791} & \num{81.04}\textsuperscript{*} \\
  \num{90000000} & \num{71277791} & \num{81.04}\textsuperscript{*} \\
  \bottomrule
\end{tabular}

\end{table}

\subsection*{L'azote}

Le second front porte sur l'azote, et il a été tracé deux fois : sous la configuration de
référence, et sous la politique agroécologique P8. La comparaison donne un résultat qu'aucun
des deux fronts ne donnerait seul : \textbf{le kilogramme d'azote est \emph{moins} cher à
retirer sous une politique déjà contraignante}. La pente moyenne vaut
\num{\frontazoteRefPente} euros par kilogramme sous la référence, contre
\num{\frontazoteHuitPente} sous P8 — parce que P8 impose simultanément un plafond d'\gls{ift},
un plafond d'eau, un plancher d'emploi et une inertie de \SI{40}{\percent} par rapport à
l'assolement existant : l'assolement qu'elle autorise a déjà renoncé aux usages les plus
intensifs, et les hectares faciles à décarboner ont déjà été pris.

Ce résultat est invisible sans tracer les fronts. Les duals locaux des deux configurations sont
quasi identiques — \num{\frontazoteRefDualBas} et \num{\frontazoteHuitDualBas} euros par
kilogramme au point le plus serré — et l'on aurait conclu à l'égalité des coûts marginaux. Un
verdict a d'ailleurs dû être renversé en cours de route : un front par $\varepsilon$-contrainte
ne se paramètre pas sur le réalisé de la configuration de référence, mais sur la borne de sa
\textbf{politique hôte}. Le coût marginal de l'azote n'est pas le même à \SI{70}{\percent} et à
\SI{30}{\percent} d'inertie, et c'est précisément ce qui rend l'objet intéressant.

Une réserve, enfin, que je ne peux pas lever. Sur le front tracé sous P8, seules
\num{\frontazoteHuitProuves} des \num{\frontazoteHuitPoints} extrémités sont prouvées optimales,
contre \num{\frontazoteRefProuves} sur \num{\frontazoteRefPoints} pour le front de référence.
Entre elles, deux points rendent \emph{exactement} le même objectif pour \emph{exactement} la
même quantité d'azote : la chaîne d'amorçage les a épinglés sur l'incumbent du point le plus
serré. \textbf{Ce palier est un artefact de solveur, pas un plateau économique.} Le front
\emph{borne} le coût marginal ; il ne le mesure pas point par point. Pour mémoire, les prix
duals des contraintes nommées, lus sur la relaxation linéaire, situent un point d'\gls{ift} à
\SI{286}{\euros}, un kilogramme d'azote à \SI{9.81}{\euros} et une heure de travail à
\SI{12.50}{\euros}.

\section{Ce que le dispositif ne permet pas de conclure}
\label{sec:limites}

Cette section n'est pas une liste d'excuses, et c'est ce qui la distingue d'une clause de
style : chaque limite y est \emph{établie par une mesure}. C'est aussi ce qui rend le reste du
chapitre citable.

\paragraph{Un solve arrêté est faux, pas imprécis.} La politique P8 non forcée a été résolue
trois fois sur le \emph{même} ensemble faisable : \SI{\mesHuitUneHeure}{\euros} après une
heure, \SI{\mesHuitTroisHeures}{\euros} après trois heures en repartant de la précédente, et
\SI{\mesHuitOptimum}{\euros} finalement prouvés optimaux. Trois réponses à une seule question.
Et la borne \gls{lp} ne mesure pas ce défaut : elle aurait annoncé le deuxième run à
\SI{\mesHuitEcartApparent}{\percent} de l'optimum quand son écart réel était de
\SI{\mesHuitEcartReel}{\percent} — cinq fois moins —, le saut d'intégralité de
\SI{\mesHuitSautInt}{\percent} en absorbant l'essentiel. \emph{Une borne \gls{lp} prouve une
infaisabilité ; elle ne chiffre pas ce qu'une recherche arborescente a laissé.} Le même
problème, avec la même graine, a par ailleurs demandé \SI{\mesHuitPreuveDuree}{\second} puis
\SI{10856}{\second} pour être prouvé : la durée elle-même n'est pas une donnée stable.

\paragraph{La même erreur s'est produite dans la grille, et sa correction a été prédite.} F0
étant le forçage \emph{nul}, le statu quo et le statu quo croisé avec F0 sont le même problème
et doivent donner le même chiffre. Ils différaient de \SI{\prosEcartFzero}{\euros}, soit
\SI{\prosEcartFzeroPct}{\percent}. La prédiction faite alors — l'écart est entièrement
imputable au hasard du branchement, une reprise depuis une graine correcte doit le ramener à
zéro — a été vérifiée au centime, en \SI{\prosEcartFzeroDuree}{\second} de calcul. Ce que cet
épisode a coûté est instructif : le classement des politiques avait tenu, mais \emph{sa
magnitude était fausse d'un facteur deux}, une cellule ayant dû être corrigée de
\SI{13.66}{\percent}.

\paragraph{Un front peut être entièrement inactif.} Sous P8, cinq des sept seuils du balayage
d'azote étaient supérieurs au plafond que la politique s'impose déjà : ils ne contraignaient
rien, et le front aurait rendu cinq fois la même allocation. Le défaut est dans le plan, pas
dans le modèle, et il ne se voit pas à la lecture du fichier — il se voit au contrôle des
points décrit au §~\ref{sec:fronts}.

\paragraph{Une politique peut être hors d'atteinte sans être infaisable.} La politique de
bifurcation P10 est faisable en relaxation \gls{lp}, dont la borne vaut
\SI{\mesPdixBorneLp}{\euros}, mais aucun solve n'en a trouvé la moindre solution entière en
\SI{7200}{\second}, et aucune graine disponible n'était à distance de réparation — la meilleure
violait encore \num{\mesPdixViolationsA} contraintes. Il en résulte que \textbf{l'axe politique
du dispositif est unilatéral} : sa borne dérégulatrice a tourné, sa borne de bifurcation non.
Ce qui reste citable de P10 est une borne supérieure, ce qui n'est pas un résultat.

\paragraph{Trois lectures sont interdites.} Premièrement, un indicateur qu'une \emph{contrainte
fixe} est une hypothèse du scénario et non son résultat : noter P8 sur son plafond d'azote
revient à la féliciter d'avoir respecté sa propre consigne. Deuxièmement, un total n'est pas une
intensité — un scénario qui cultive moins d'hectares réduit l'azote total sans rien changer aux
pratiques. Troisièmement, le \gls{pad} du chapitre~\ref{ch:calibrer} ne s'applique pas ici : un
scénario prospectif est construit pour s'écarter de 2017, et le générateur de chiffres de ce
mémoire refuse d'ailleurs d'émettre un \gls{pad} pour un run prospectif. À quoi s'ajoutent deux
indicateurs à manier avec précaution : un ratio d'autonomie alimentaire suppose une
\emph{demande}, et la demande n'est pas modélisée ; et une part minimale de cultures
biologiques bâtie sur les vingt-cinq variantes symétriques du §~\ref{sec:tractabilite} est
satisfaite \textbf{à coût exactement nul}, par relabellisation.

\paragraph{Et la réserve la plus lourde, qui conditionne le résultat le plus spectaculaire.}
Aucun débouché n'est modélisé au-delà des deux filières d'exportation. Le basculement de
\SI{2100}{\ha} vers le vivrier décrit au §~\ref{sec:fronts} repose donc sur une hypothèse de
marché \emph{parfaitement élastique} pour le maraîchage, les agrumes et l'igname. Le modèle dit
ce que l'agronomie et l'économie de la \emph{production} rendent possible ; il ne dit pas si le
territoire absorberait cette production, ni à quel prix. C'est très exactement l'information que
les enquêtes prévues par la lettre de mission devaient apporter.

\begin{reflexion}
Le catalogue de dix politiques publiques que ce chapitre exploite, je l'ai écrit seul :
l'axe, les seuils, et le choix de faire du statu quo P4 le point de référence implicite de
toute comparaison. Ce sont des décisions politiques, prises par un stagiaire dans un fichier
\gls{yaml}, sur la base de sa lecture de la littérature et de rien d'autre. Ce qui empêche que
ce soit une faute est qu'il faut bien que quelqu'un écrive le premier catalogue, et que le
rendre explicite et modifiable ligne à ligne est précisément ce qui permet à quelqu'un d'autre
de le contester. Mais l'argument ne suffit pas. L'axe qui organise le catalogue — de la
dérégulation à la bifurcation — suppose que le débat agricole guadeloupéen s'ordonne ainsi :
c'est une hypothèse, pas une observation. Et l'ancrage mesuré est cruel, puisque seule la borne
dérégulatrice a effectivement tourné. Le projet qui a précédé celui-ci avait mené une démarche
participative — séminaire, \emph{World Café}, questionnaire ; la conception participative des
scénarios n'aurait pas été un supplément d'âme, elle aurait changé les seuils, très probablement
les leviers, peut-être l'axe lui-même.\attente{3} Ce que le dispositif produit doit donc être lu
pour ce qu'il est : les conséquences rigoureusement calculées d'hypothèses que personne du
territoire n'a validées. \textbf{La rigueur du calcul ne se transmet pas aux hypothèses.}
\end{reflexion}
